% !TeX root = main_long.tex
\section{Discussion}
\label{sec:discussion}

% Flow: coverage and findings; gaps; converter assessment; integration utility;
% position among models; future work and independent review.

\subsection{Assessing VCF Core's coverage of VCF meaning}
\label{sec:summary}

This study assesses whether VCF Core makes the information needed to interpret VCF observations available for reuse in RDF.
Preserving a field as text does not establish that its declaration, allele ordering, and sample context are accessible to a graph query.
RQ1 therefore tests preservation and structured retrieval, while RQ2 examines whether the same answers can be obtained from a separately implemented converter among other assessments designed to test the robustness of \textit{vcfcv}'s stated goal of faithfully representing the VCF specificaiton.
The results support broad representation with incomplete demonstrated coverage: the available expanded-profile tests pass, but many requirements remain unassessed (Table~\ref{tab:requirement-results}).
The lower condensed-structure scores reflect a different access mechanism, since individual sample values require decoding.
For users, this distinction makes coverage meaningful in relation to both the information needed and the way an application must retrieve it.

\subsection{Extending demonstrated coverage}

The clearest opportunity to improve coverage is to extend testing of structural-variant fields, allele-dependent values, and version-specific rules.
New fixtures should isolate each rule, establish expected answers from the specification, and exercise all applicable versions~\cite{samod2016,blomqvist2012ontologytesting}.
Validation also needs to extend to constructs without enforcement evidence.
Because clear specification examples guided the current tests, the demonstrated subset should not be treated as representative of all VCF complexity.

Deriving executable questions also exposed inconsistencies in specification examples and declarations.
This makes vocabulary assessment useful for specification maintenance: an executable example gives maintainers a concrete interpretation to review.
Similar benefits have been reported for protocol specifications and genomics-format conformance~\cite{bishop2005conformance,niu2022interop}, while unstated VCF conventions have also complicated reuse~\cite{beier2022fairvcf}.
Clarifying these rules, refining generic requirement descriptions, and testing files from different production pipelines would strengthen the evidence beyond the present fixtures.

\subsection{A reusable assessment for conversion engines}

RQ2 shows how the same assessment can guide converter development.
The corrected local-allele error produced plausible values attached to the wrong alleles, while the remaining empty-list mismatch was detected through complementary query and round-trip checks.
These findings demonstrate why validation of graph structure should be accompanied by checks of recovered values and their associations.

A shared vocabulary makes those checks reusable across conversion strategies.
Implementers can develop streaming, mapping-based, or materializing converters against the same representation, while users can evaluate the fields and sample profile their application requires before comparing runtime or graph size.
The current comparison supports this approach across two separate implementations, although their shared authorship leaves open the possibility of a common misunderstanding.
Re-execution by unrelated groups is therefore a necessary next step.

\subsection{Integrating observations with external knowledge}

We chose local alleles for RQ3 because they make the importance of interpretation context particularly clear.
A record can list alternate alleles observed across a cohort, while a sample's \texttt{LAA} identifies a subset.
Depths in \texttt{LAD}, declared with \texttt{Number=LR}, refer to the reference allele and that local subset.
Retrieving depth for an externally annotated alteration therefore depends on the declaration, the sample mapping, and the record allele together.
The specification recommends abstracting this encoding within libraries~\cite{vcf45}; VCF Core exposes the resolved association through \term{forAllele}.

The preliminary findings are encouraging because this relationship changes which alterations satisfy the query, even when both interpretations yield plausible rows with provenance.
Its practical value is that the interpretation can be implemented and checked during conversion, then reused by downstream queries without decoding local alleles again.
The synthetic example assumes a shared assembly and identical written alleles; it demonstrates this capability while leaving realistic biomedical integration and reference-aware normalization to further evaluation.

\subsection{VCF Core among existing models}

RQ4 identifies VCF Core's distinct role in making source interpretation context directly addressable.
The inspected models serve complementary purposes, while some VCF relationships require parsing or additional conventions (Table~\ref{tab:models}).
The criteria reflect the problems VCF Core was designed to address, so the comparison should be read within that scope.

Observations from two laboratories could retain their own quality values, sample context, and provenance while linking, after normalization, to a common variation description such as GA4GH VRS~\cite{wagner2021vrs}.
A clinical model such as HERO-Genomics could connect sample columns to specimens and patients~\cite{cazzaro2025hero}; Med2RDF and \emph{gvar} provide further candidates for alignment~\cite{med2rdf,gvar}.
The published mappings provide graded relationships between terms, rather than identity between instances.
Review by the respective model communities would strengthen these connections and the inspection-based comparison~\cite{smith2007obo}.

\subsection{Future work and independent evaluation}
\label{sec:clinicalqueries}

Future studies should test this foundation in workflows linking pharmacogenomic interpretations to supporting observations~\cite{lee2022cpic,pharmcatmethods}, rare-disease candidates to phenotype knowledge~\cite{hpo,putman2024monarch}, or tumour observations to therapeutic evidence~\cite{griffith2017civic}.
These proposed scenarios require biomedical models, reviewed instance mappings, and explicit sample--specimen--patient relationships.
Their value should be assessed through answer correctness and traceability.

Realistic cohort studies should also measure conversion, storage, loading, and query costs alongside correctness.
Current round-trip and integration results apply only to expanded data; SPARQL extensions such as SciSPARQL~\cite{andrejev2012scisparql} could support investigation of access to condensed vectors.
We invite independent review of the vocabulary, converters, requirement interpretations, and expected answers.
New fixtures, counterexamples, and disputed interpretations would directly improve this resource and its basis in the VCF specification.
